\documentclass[../../main.tex]{subfiles}
\graphicspath{{\subfix{../../image/}}}

\begin{document}

\section{完全积分、包络}\label{section:ch03-complete}

我们从分析非线性一阶 PDE
\begin{align}\label{align:ch03-pde}
F(Du,u,x)=0
\end{align}
入手：先描述一些简单的解类，再学习如何由它们构造更复杂的解。

先设 $A\subset\mathbb{R}^n$ 为开集。假定对每个参数 $a=(a_1,\cdots,a_n)\in A$，我们都有 PDE \eqref{align:ch03-pde} 的一个 $C^2$ 解 $u=u(x;a)$。

\textbf{记号.} 我们写
\begin{align}\label{align:ch03-dau}
(D_a u,D_{xa}^2 u)\triangleq
\begin{pmatrix}
u_{a_1} & u_{x_1 a_1} & \cdots & u_{x_n a_1}\\
\vdots & \vdots & \ddots & \vdots\\
u_{a_n} & u_{x_1 a_n} & \cdots & u_{x_n a_n}
\end{pmatrix}_{n\times(n+1)}.
\end{align}

\begin{definition}\label{definition:ch03-complete}
一个 $C^2$ 函数 $u=u(x;a)$ 称为在 $U\times A$ 中的\textbf{完全积分}，只要
\begin{enumerate}[(i)]
\item 对每个 $a\in A$，$u(x;a)$ 解 PDE \eqref{align:ch03-pde}；
\item $\operatorname{rank}(D_a u,D_{xa}^2 u)=n$（$x\in U,\ a\in A$）。
\end{enumerate}
\end{definition}

\textbf{解释.} 条件 (ii) 保证 $u(x;a)$``依赖于全部 $n$ 个独立参数 $a_1,\cdots,a_n$''。为看出此点，设 $B\subset\mathbb{R}^{n-1}$ 为开集，对每个 $b\in B$ 假定 $v=v(x;b)$（$x\in U$）是\eqref{align:ch03-pde} 的解。又设存在 $C^1$ 映射 $\boldsymbol{\psi}:A\to B$，$\boldsymbol{\psi}=(\psi_1,\cdots,\psi_{n-1})$，使
\begin{align}\label{align:ch03-depends-n-1}
u(x;a)=v(x;\boldsymbol{\psi}(a)) \qquad (x\in U,\ a\in A).
\end{align}
即我们假定函数 $u(x;a)$``实际上只依赖于 $n-1$ 个参数 $b_1,\cdots,b_{n-1}$''。但此时
\begin{align*}
u_{x_i a_j}(x;a)=\sum_{k=1}^{n-1}v_{x_i b_k}(x;\boldsymbol{\psi}(a))\frac{\partial\psi_k}{\partial a_j}(a) \qquad (i,j=1,\cdots,n).
\end{align*}
因而
\begin{align*}
\det(D_{xa}^2 u)=\sum_{k_1,\cdots,k_n=1}^{n-1}v_{x_1 b_{k_1}}\cdots v_{x_n b_{k_n}}\det
\begin{pmatrix}
\dfrac{\partial\psi_{k_1}}{\partial a_1} & \cdots & \dfrac{\partial\psi_{k_1}}{\partial a_n}\\
& \ddots & \\
\dfrac{\partial\psi_{k_n}}{\partial a_1} & \cdots & \dfrac{\partial\psi_{k_n}}{\partial a_n}
\end{pmatrix}=0,
\end{align*}
因为对 $k_1,\cdots,k_n\in\{1,\cdots,n-1\}$ 的每一种选取，相应矩阵中至少两行相同。又
\begin{align*}
u_{a_j}(x;a)=\sum_{k=1}^{n-1}v_{b_k}(x;\boldsymbol{\psi}(a))\frac{\partial\psi_k}{\partial a_j}(a) \qquad (j=1,\cdots,n),
\end{align*}
类似论证表明 $(D_a u,D_{xa}^2 u)$ 的每个 $n\times n$ 子矩阵的行列式都等于零，故该矩阵秩严格小于 $n$。

\begin{example}\label{example:ch03-clairaut}
微分几何中的 Clairaut 方程是 PDE
\begin{align}\label{align:ch03-clairaut}
x\cdot Du+f(Du)=u,
\end{align}
其中 $f:\mathbb{R}^n\to\mathbb{R}$ 给定。一个完全积分为
\begin{align}\label{align:ch03-clairaut-int}
u(x;a)=a\cdot x+f(a) \qquad (x\in U)
\end{align}
对 $a\in\mathbb{R}^n$。
\end{example}

\begin{example}\label{example:ch03-eikonal}
几何光学中的 eikonal$^*$ 方程是 PDE
\begin{align}\label{align:ch03-eikonal}
|Du|=1.
\end{align}
一个完全积分为
\begin{align}\label{align:ch03-eikonal-int}
u(x;a,b)=a\cdot x+b \qquad (x\in U)
\end{align}
其中 $x\in U$，$a\in\partial B(0,1)$，$b\in\mathbb{R}$。

\smallskip
\noindent $^*$ $\varepsilon\acute{\iota}\kappa\acute{\omega}\nu=$ image（希腊语）。
\end{example}

\begin{example}\label{example:ch04-3-3-8}
力学中的 Hamilton--Jacobi 方程最简形式为 PDE
\begin{align}\label{align:ch03-hj}
u_t+H(Du)=0,
\end{align}
其中 $H:\mathbb{R}^n\to\mathbb{R}$。这里 $u$ 依赖于 $x=(x_1,\cdots,x_n)\in\mathbb{R}^n$ 与 $t\in\mathbb{R}$。如前令 $t=x_{n+1}$，写 $Du=D_x u=(u_{x_1},\cdots,u_{x_n})$。一个完全积分为
\begin{align}\label{align:ch03-hj-int}
u(x,t;a,b)=a\cdot x-tH(a)+b \qquad (x\in\mathbb{R}^n,\ t\geqslant 0)
\end{align}
其中 $a\in\mathbb{R}^n$，$b\in\mathbb{R}$。
\end{example}

\subsection{由包络构造新解}\label{subsection:ch03-envelope}

接下来我们说明如何构造非线性一阶 PDE \eqref{align:ch03-pde} 的更复杂解——这些解依赖于 $n-1$ 个变量的任意函数，而不仅仅依赖于 $n$ 个参数。我们将把这些新解构造为完全积分（或更一般地，其他 $m$ 参数解族）的包络。

\begin{definition}\label{definition:ch03-envelope}
设 $u=u(x;a)$ 是 $x\in U$、$a\in A$ 的 $C^1$ 函数，其中 $U\subset\mathbb{R}^n$、$A\subset\mathbb{R}^m$ 为开集。考虑向量方程
\begin{align}\label{align:ch03-da-u=0}
D_a u(x;a)=0 \qquad (x\in U,\ a\in A).
\end{align}
设我们能把\eqref{align:ch03-da-u=0} 中的参数 $a$ 解出为 $x$ 的 $C^1$ 函数：
\begin{align}\label{align:ch03-a=phi}
a=\boldsymbol{\phi}(x),
\end{align}
从而
\begin{align}\label{align:ch03-da-u=0-at-phi}
D_a u(x;\boldsymbol{\phi}(x))=0 \qquad (x\in U).
\end{align}
我们称
\begin{align}\label{align:ch03-envelope-def}
v(x)\triangleq u(x;\boldsymbol{\phi}(x)) \qquad (x\in U)
\end{align}
为函数族 $\{u(\cdot;a)\}_{a\in A}$ 的\textbf{包络}。
\end{definition}

通过作包络，我们可以构造非线性一阶 PDE 的新解：

\begin{theorem}[新解的构造]\label{theorem:ch03-envelope}设如上每个 $a\in A$ 对应的 $u=u(\cdot;a)$ 解 PDE \eqref{align:ch03-pde}。又设由上面\eqref{align:ch03-da-u=0-at-phi}、\eqref{align:ch03-envelope-def} 定义的包络 $v$ 存在且为 $C^1$ 函数。则 $v$ 也解\eqref{align:ch03-pde}。
\end{theorem}

上面定义的包络 $v$ 有时称为\eqref{align:ch03-pde} 的\textbf{奇异积分}。

\begin{proof}
我们有 $v(x)=u(x;\boldsymbol{\phi}(x))$；故对 $i=1,\cdots,n$，
\begin{align*}
v_{x_i}(x)=u_{x_i}(x;\boldsymbol{\phi}(x))+\sum_{j=1}^m u_{a_j}(x,\boldsymbol{\phi}(x))\frac{\partial\phi_j}{\partial x_i}(x)
=u_{x_i}(x;\boldsymbol{\phi}(x)),
\end{align*}
最后一个等号由\eqref{align:ch03-da-u=0-at-phi} 得到。因而对每个 $x\in U$，
\begin{align*}
F(Dv(x),v(x),x)=F(Du(x;\boldsymbol{\phi}(x)),u(x;\boldsymbol{\phi}(x)),x)=0.
\end{align*}
\end{proof}

几何意义是：对每个 $x\in U$，$v$ 的图像在 $a=\boldsymbol{\phi}(x)$ 处与 $u(\cdot;a)$ 的图像相切。因而在 $x$ 处（$a=\boldsymbol{\phi}(x)$）有 $Dv=D_x u(\cdot;a)$。

\begin{example}\label{example:ch03-sphere}
考虑 PDE
\begin{align}\label{align:ch03-sphere-pde}
u^2(1+|Du|^2)=1.
\end{align}
一个完全积分为 $u(x,a)=\pm(1-|x-a|^2)^{1/2}$（$|x-a|<1$）。计算得
\begin{align*}
D_a u=\mp\frac{x-a}{(1-|x-a|^2)^{1/2}}=0
\end{align*}
只要 $a=\boldsymbol{\phi}(x)=x$。因而 $v\equiv\pm 1$ 是\eqref{align:ch03-sphere-pde} 的奇异积分。
\end{example}

为从一个完全积分生成 PDE \eqref{align:ch03-pde} 的更多解，我们改造上述构造。任取开集 $A'\subset\mathbb{R}^{n-1}$ 与任一 $C^1$ 函数 $h:A'\to\mathbb{R}$，使 $h$ 的图像落在 $A$ 内。写 $a=(a_1,\cdots,a_n)=(a',a_n)$，其中 $a'=(a_1,\cdots,a_{n-1})\in\mathbb{R}^{n-1}$。

\begin{definition}[依赖于 $h$ 的]\label{definition:ch03-general}\textbf{一般积分}是函数族
\begin{align*}
u'(x;a')=u(x;a',h(a')) \qquad (x\in U,\ a'\in A')
\end{align*}
的包络 $v'=v'(x)$，只要此包络存在且为 $C^1$。
\end{definition}

换言之，在计算包络时，我们现在只限制于形如 $a=(a',h(a'))$ 的参数 $a$，其中函数 $h$ 作某一显式选取。于是从一个依赖于 $n$ 个任意常数 $a_1,\cdots,a_n$ 的完全积分，我们（只要上述构造可行）构造出一个依赖于 $n-1$ 个变量的任意函数 $h$ 的解。

\begin{example}\label{example:ch03-hj-envelope}
在上面\cref{example:ch04-3-3-8} 中令 $H(p)=|p|^2$，$h\equiv 0$。则
\begin{align*}
u'(x,t;a)=x\cdot a-t|a|^2.
\end{align*}
令 $D_a u'=x-2ta=0$ 计算包络。故 $a=\dfrac{x}{2t}$，从而
\begin{align*}
v'(x,t)=x\cdot\frac{x}{2t}-t\left|\frac{x}{2t}\right|^2=\frac{|x|^2}{4t} \qquad (x\in\mathbb{R}^n,\ t>0)
\end{align*}
解 Hamilton--Jacobi 方程 $v'_t+|Dv'|^2=0$。
\end{example}

\begin{remark}\label{remark:ch03-envelope-all}
人们很容易相信：一旦如上找到一个依赖于任意函数 $h$ 的\eqref{align:ch03-pde} 的解，便找到了\eqref{align:ch03-pde} 的所有解。但这未必成立。设我们的 PDE 具有结构
\begin{align*}
F(Du,u,x)=F_1(Du,u,x)F_2(Du,u,x)=0.
\end{align*}
若 $u_1(x,a)$ 是 PDE $F_1(Du,u,x)=0$ 的一个完全积分，且我们成功找到对应于任一函数 $h$ 的一般积分，我们仍会漏掉 PDE $F_2(Du,u,x)=0$ 的所有解。
\end{remark}

\end{document}